\section{Ensemblové modelování v regresních a klasifikačních úlohách. Bagging a boosting}

\subsection*{Základní idea ensemblů}

Ensemble model kombinuje více jednodušších modelů do jednoho výsledného modelu. Místo jednoho klasifikátoru
nebo regresoru použijeme skupinu modelů:
\[
f_1, f_2, \ldots, f_M.
\]

Výsledná predikce vznikne agregací.

Pro regresi často průměrováním:
\[
\hat{y}(x)
=
\frac{1}{M}
\sum_{m=1}^M
f_m(x).
\]

Pro klasifikaci hlasováním:
\[
\hat{y}(x)
=
\operatorname{mode}
\{f_1(x),\ldots,f_M(x)\}.
\]

\paragraph{Intuice.}
Jeden model může být náhodně vychýlený nebo citlivý na konkrétní trénovací data. Pokud zkombinujeme více
různých modelů, chyby jednotlivých modelů se mohou částečně vyrušit.

\subsection*{Proč ensemble metody fungují}

Ensemble metody pomáhají hlavně tehdy, když základní modely:
\begin{itemize}
    \item jsou lepší než náhoda,
    \item nejsou úplně stejné,
    \item dělají různé chyby.
\end{itemize}

Diverzita modelů je klíčová. Pokud by všechny modely byly totožné, ensemble by nepřinesl žádné zlepšení.

\subsection*{Bias-variance intuice}

Chybu modelu lze intuitivně rozložit na:
\[
\text{chyba}
=
\text{bias}
+
\text{variance}
+
\text{noise}.
\]

\begin{itemize}
    \item \textbf{Bias} je chyba způsobená příliš jednoduchým modelem.
    \item \textbf{Variance} je citlivost modelu na konkrétní trénovací data.
    \item \textbf{Noise} je neodstranitelný šum v datech.
\end{itemize}

Bagging typicky snižuje hlavně varianci. Boosting typicky snižuje bias, ale při špatném nastavení může zvýšit
riziko přeučení.

\subsection*{Bagging}

Bagging znamená \term{bootstrap aggregating}. Vytváříme více trénovacích množin pomocí výběru s navracením.

\subsubsection*{Bootstrap sample}

Z původních dat:
\[
D=\{(x_i,y_i)\}_{i=1}^n
\]
vytvoříme bootstrapový vzorek $D_m$ tak, že náhodně vybereme $n$ pozorování s navracením.

Jedno pozorování se tedy může objevit vícekrát a jiné vůbec.

\subsubsection*{Postup baggingu}

\begin{enumerate}
    \item vytvoř $M$ bootstrapových vzorků $D_1,\ldots,D_M$,
    \item na každém vzorku natrénuj model $f_m$,
    \item pro regresi průměruj predikce,
    \item pro klasifikaci nech modely hlasovat.
\end{enumerate}

Regrese:
\[
\hat{y}(x)=
\frac{1}{M}
\sum_{m=1}^M
f_m(x).
\]

Klasifikace:
\[
\hat{y}(x)
=
\underset{k}{\operatorname{argmax}}
\sum_{m=1}^M
I(f_m(x)=k).
\]

\subsection*{Kdy bagging pomáhá}

Bagging je vhodný hlavně pro nestabilní modely, kde malá změna trénovacích dat může výrazně změnit model.
Typický příklad jsou rozhodovací stromy.

Bagging pomáhá:
\begin{itemize}
    \item snížit varianci,
    \item omezit přeučení,
    \item zlepšit robustnost,
    \item stabilizovat predikce.
\end{itemize}

\subsection*{Out-of-bag error}

Při bootstrapu se některá pozorování do daného bootstrapového vzorku nedostanou. Tato pozorování jsou
\term{out-of-bag} pro daný model.

OOB pozorování lze využít pro odhad testovací chyby bez samostatné validační množiny.

\[
Err_{OOB}
=
\frac{1}{n}
\sum_{i=1}^n
L(y_i,\hat{y}^{OOB}_i).
\]

\subsection*{Random subspaces, pasting a random patches}

Materiály uvádějí i varianty baggingu:

\begin{itemize}
    \item \textbf{Pasting:}
    vybírají se náhodné příklady bez navracení.
    
    \item \textbf{Random subspaces:}
    pro každý model se vybírá náhodná podmnožina atributů.
    
    \item \textbf{Random patches:}
    kombinuje se výběr náhodných příkladů i náhodných atributů.
\end{itemize}

\subsection*{Random forest}

Random forest je ensemble rozhodovacích stromů. Používá dvě hlavní formy náhodnosti:

\begin{enumerate}
    \item bootstrapový výběr pozorování,
    \item náhodný výběr části atributů při hledání nejlepšího splitu.
\end{enumerate}

Predikce v klasifikaci:
\[
\hat{y}(x)
=
\underset{k}{\operatorname{argmax}}
\sum_{m=1}^M
I(T_m(x)=k).
\]

Predikce v regresi:
\[
\hat{y}(x)
=
\frac{1}{M}
\sum_{m=1}^M
T_m(x).
\]

\paragraph{Proč nestačí jen bagging stromů?}
Rozhodovací stromy mohou být mezi sebou silně korelované, hlavně pokud existuje několik velmi silných
prediktorů. Random forest snižuje korelaci stromů tím, že při každém splitu zvažuje jen náhodnou podmnožinu
atributů.

\subsection*{Důležité hyperparametry random forest}

\begin{itemize}
    \item \texttt{n\_estimators} -- počet stromů,
    \item \texttt{max\_features} -- počet atributů zvažovaných při splitu,
    \item \texttt{max\_depth} -- maximální hloubka stromu,
    \item \texttt{min\_samples\_leaf} -- minimální počet pozorování v listu,
    \item \texttt{bootstrap} -- zda používat bootstrap,
    \item \texttt{max\_samples} -- velikost bootstrapového vzorku.
\end{itemize}

\subsection*{Výhody a nevýhody random forest}

Výhody:
\begin{itemize}
    \item vysoká predikční přesnost,
    \item robustnost vůči šumu,
    \item menší sklon k přeučení než jeden strom,
    \item možnost měřit důležitost proměnných,
    \item funguje pro klasifikaci i regresi.
\end{itemize}

Nevýhody:
\begin{itemize}
    \item nižší interpretovatelnost než jeden strom,
    \item větší výpočetní náročnost,
    \item u velmi nevyvážených dat je nutné řešit váhy tříd nebo sampling,
    \item importance proměnných může být zkreslená.
\end{itemize}

\subsection*{Boosting}

Boosting je ensemble metoda, která modely staví sekvenčně. Na rozdíl od baggingu nejsou modely nezávislé.

Obecný tvar:
\[
\hat{y}_i
=
f(x_i)
=
\sum_{m=1}^M
f_m(x_i;\Theta_m).
\]

Každý další model se snaží opravit chyby předchozích modelů.

\paragraph{Intuice.}
Bagging trénuje modely paralelně a potom je průměruje. Boosting trénuje modely postupně: první model udělá
hrubou predikci, další model se učí na tom, co první nezvládl, další opravuje zbylé chyby atd.

\subsection*{Weak learner}

Boosting kombinuje mnoho slabých modelů do silného modelu.

Slabý model je model, který je jen o něco lepší než náhodné hádání. Typicky se používají malé rozhodovací
stromy, například decision stumps:
\[
\text{strom s hloubkou } 1.
\]

\subsection*{Forward stagewise additive fitting}

Boosting lze chápat jako aditivní model:
\[
F_M(x)
=
F_0(x)
+
\sum_{m=1}^M
\nu f_m(x),
\]
kde $\nu$ je learning rate.

Inicializace:
\[
F_0(x)=\bar{y}
\]
u regrese s kvadratickou ztrátou.

V kroku $m$ se hledá nový model:
\[
\Theta_m
=
\underset{\Theta}{\operatorname{argmin}}
\sum_{i=1}^n
L
\left(
y_i,
F_{m-1}(x_i)+f_m(x_i;\Theta)
\right).
\]

Aktualizace:
\[
F_m(x_i)
=
F_{m-1}(x_i)
+
\nu f_m(x_i;\Theta_m).
\]

\subsection*{Boosting u kvadratické ztráty}

Pro regresi s kvadratickou ztrátou:
\[
L(y,\hat{y})=(y-\hat{y})^2.
\]

Rezidua po kroku $m-1$:
\[
r_{i,m-1}
=
y_i-F_{m-1}(x_i).
\]

Nový strom se učí predikovat rezidua:
\[
\Theta_m
=
\underset{\Theta}{\operatorname{argmin}}
\sum_{i=1}^n
\left(
r_{i,m-1}-f_m(x_i;\Theta)
\right)^2.
\]

Potom:
\[
F_m(x_i)
=
F_{m-1}(x_i)
+
\nu f_m(x_i;\Theta_m).
\]

\paragraph{Intuice.}
Pokud aktuální model u některých pozorování predikuje moc nízko, reziduum je kladné a další strom se je snaží
zvýšit. Pokud predikuje moc vysoko, reziduum je záporné a další strom ho snižuje.

\subsection*{Gradient boosting}

Gradient boosting zobecňuje boosting na libovolnou diferencovatelnou ztrátovou funkci.

Místo obyčejných reziduí se používá záporný gradient ztrátové funkce:
\[
r_{im}
=
-
\left[
\frac{\partial L(y_i,F(x_i))}
{\partial F(x_i)}
\right]_{F=F_{m-1}}.
\]

Nový slabý model se učí aproximovat tento záporný gradient:
\[
f_m
\approx
r_{im}.
\]

Aktualizace:
\[
F_m(x)
=
F_{m-1}(x)
+
\nu f_m(x).
\]

\paragraph{Intuice.}
Gradient ukazuje směr nejrychlejšího růstu ztráty. Záporný gradient ukazuje směr, kterým musíme predikce
změnit, aby ztráta klesala.

\subsection*{Learning rate}

Learning rate $\nu$ určuje velikost kroku:
\[
F_m(x)=F_{m-1}(x)+\nu f_m(x).
\]

Malá hodnota $\nu$:
\begin{itemize}
    \item zpomaluje učení,
    \item obvykle zlepšuje generalizaci,
    \item vyžaduje více stromů.
\end{itemize}

Velká hodnota $\nu$:
\begin{itemize}
    \item rychle snižuje trénovací chybu,
    \item zvyšuje riziko přeučení.
\end{itemize}

\subsection*{AdaBoost}

AdaBoost je klasický boostingový algoritmus pro klasifikaci. Myšlenka je měnit váhy pozorování.

Po každém modelu:
\begin{itemize}
    \item špatně klasifikovaným pozorováním se zvýší váha,
    \item správně klasifikovaným pozorováním se váha sníží.
\end{itemize}

Chyba modelu:
\[
err_m
=
\frac{
\sum_{i=1}^n w_i I(y_i\neq f_m(x_i))
}{
\sum_{i=1}^n w_i
}.
\]

Váha modelu:
\[
\alpha_m
=
\frac{1}{2}
\ln
\left(
\frac{1-err_m}{err_m}
\right).
\]

Finální klasifikátor:
\[
F(x)
=
\operatorname{sign}
\left(
\sum_{m=1}^M
\alpha_m f_m(x)
\right).
\]

\subsection*{XGBoost}

XGBoost je efektivní implementace gradient boostingu nad stromy. Je populární hlavně pro tabulková data.

Používá regularizovanou ztrátovou funkci:
\[
Obj^{(m)}
=
\sum_{i=1}^n
L
\left(
y_i,
\hat{y}^{(m-1)}_i
+
f_m(x_i)
\right)
+
\Omega(f_m).
\]

Regularizace stromu:
\[
\Omega(f)
=
\gamma J
+
\frac{1}{2}
\lambda
\sum_{j=1}^J
c_j^2.
\]

Kde:
\begin{itemize}
    \item $J$ je počet listů,
    \item $c_j$ je predikce v listu $j$,
    \item $\gamma$ penalizuje počet listů,
    \item $\lambda$ penalizuje velikost predikcí v listech.
\end{itemize}

Cílem je nejen snížit chybu, ale i omezit složitost stromů.

\subsection*{Důležité hyperparametry boostingu}

\begin{itemize}
    \item počet stromů $M$,
    \item learning rate $\nu$,
    \item maximální hloubka stromů,
    \item minimální počet pozorování v listu,
    \item subsampling pozorování,
    \item subsampling atributů,
    \item regularizační parametry,
    \item early stopping.
\end{itemize}

\subsection*{Early stopping}

Boosting může být citlivý na šum a extrémní hodnoty. Proto se sleduje chyba na validačních datech.

\[
Err_{\text{val}}(m).
\]

Pokud se validační chyba po určitou dobu nezlepšuje, trénování se zastaví.

\subsection*{Bagging vs. boosting}

\begin{itemize}
    \item \textbf{Bagging:}
    modely jsou trénovány paralelně a nezávisle.
    
    \item \textbf{Boosting:}
    modely jsou trénovány sekvenčně a závisle.
    
    \item \textbf{Bagging:}
    hlavně snižuje varianci.
    
    \item \textbf{Boosting:}
    hlavně snižuje bias, ale může se přeučit.
    
    \item \textbf{Bagging:}
    používá bootstrap a hlasování/průměrování.
    
    \item \textbf{Boosting:}
    postupně opravuje chyby předchozího modelu.
    
    \item \textbf{Random forest:}
    typický zástupce baggingu.
    
    \item \textbf{Gradient boosting / XGBoost:}
    typičtí zástupci boostingu.
\end{itemize}

\subsection*{Stacking}

Materiály zmiňují také stacking. Ten kombinuje různé modely pomocí meta-modelu.

Máme modely:
\[
f_1(x), f_2(x), \ldots, f_M(x).
\]

Jejich predikce se použijí jako vstupy do meta-modelu:
\[
\hat{y}
=
g(f_1(x),f_2(x),\ldots,f_M(x)).
\]

Meta-model se učí, kterému základnímu modelu a v jaké situaci má věřit.

\subsection*{Hodnocení ensemble modelů}

Ensemble modely je nutné hodnotit na validačních nebo testovacích datech. U klasifikace:
\[
Accuracy,\quad Precision,\quad Recall,\quad F_1,\quad AUC.
\]

U regrese:
\[
MSE=
\frac{1}{n}
\sum_{i=1}^n
(y_i-\hat{y}_i)^2,
\]
\[
RMSE=
\sqrt{MSE},
\]
\[
MAE=
\frac{1}{n}
\sum_{i=1}^n
|y_i-\hat{y}_i|.
\]

\subsection*{Interpretovatelnost ensemblů}

Ensemble modely mají obvykle vyšší přesnost než jednotlivé stromy, ale nižší interpretovatelnost.

Používají se:
\begin{itemize}
    \item feature importance,
    \item permutation importance,
    \item partial dependence plots,
    \item ICE plots,
    \item SHAP hodnoty,
    \item LIME.
\end{itemize}

\paragraph{Praktický kompromis.}
Pokud je cílem auditovatelnost a vysvětlitelnost, může být vhodnější menší strom nebo pravidlový model. Pokud
je cílem predikční výkon, často vítězí random forest nebo gradient boosting.

\subsection*{Shrnutí otázky}

Ensemblové modely kombinují více modelů do jednoho výsledného modelu. Bagging vytváří více bootstrapových
vzorků, na nich trénuje nezávislé modely a jejich predikce agreguje. Typickým příkladem je random forest,
který kombinuje bootstrap pozorování s náhodným výběrem atributů.

Boosting modely trénuje sekvenčně. Každý další model opravuje chyby předchozích. Gradient boosting chápe
trénování jako postupné snižování ztrátové funkce ve směru záporného gradientu. XGBoost je efektivní
regularizovaná implementace gradient boostingu, často velmi silná pro tabulková data. Bagging je robustnější
a méně citlivý, boosting může dosáhnout vyšší přesnosti, ale vyžaduje pečlivé ladění.

\newpage